\section{The source first-minimum estimate at a half maximum}
\label{sec:sr-source-half-maximum}
We keep the source frame datum, weighted norms, and rational
transvections throughout. The first-minimum estimate is made at
one actual frame of the same uniformization used by the theta reading.

\subsection{The periodic Schwarz estimate}
\begin{srlemma}\label{lem:sr-source-periodic-schwarz}
Let \(M>0\), \(0<t<T\), and let \(g:\mathbb C\to\mathbb C\) be
\(M\)-periodic, holomorphic on \(\operatorname{Im}z>-T\), and tend
to zero as \(\operatorname{Im}z\to+\infty\), uniformly in the real
coordinate in the sense of the source filter
\(\operatorname{comap}(\operatorname{Im})(+\infty)\).
If \(|g(z)|\le B\) on \(\operatorname{Im}z=-t\), then
\[
                    |g(0)|\le e^{-2\pi t/M}B .
 \tag{HM1}\label{eq:sr-source-periodic-schwarz}
\]
\end{srlemma}
\begin{proof}
Use \(q=e^{2\pi iz/M}\). Periodicity gives a single-valued
holomorphic function \(F(q)\) on the punctured disc of radius
\(e^{2\pi T/M}\).
The given limit makes \(F\) bounded near zero with \(F(0)=0\);
the removable singularity theorem, or its Cauchy coefficient proof,
makes this a holomorphic extension.
Set \(R=e^{2\pi t/M}>1\).
The divided difference \(F(q)/q\), with value \(F'(0)\) at zero,
is holomorphic on a neighborhood of the closed disc of radius \(R\).
On its boundary it has norm at most \(B/R\).
The maximum principle gives \(|F(1)|\le B/R\), which is (HM1).
This is precisely
\srcid{PeriodicSchwarzLemma.periodic_schwarz}.
\end{proof}

\subsection{Decay on the rational isotropic flow}
The source \(X=\mathrm{DataX}(E,r)\) includes the lattice \(L'\),
the scalar \(k_0>0\) with
\[
                     c_\tau=2|\tau(\xi_0)|=k_0,
\]
and an integer \(n_0>0\) sending \(L'\) into the lattice of
admissible rational isotropic vectors.
For \(v\in L'\) with \(v\ne0\) and
\(\langle v,v\rangle=0\), let \(x(\zeta)\) be the source flow
\(X.\mathrm{flow}(v,(\zeta)_\tau,x_0)\).

We use the exact source Fourier expansion on an admissible
isotropic direction and its coefficient-positivity statement.
The latter is the first Fourier input in
Section~\ref{sec:sr-source-integer-reading}.
For \(P\in SS_\lambda(X)\), write \(u=n_0v\).
The square in the rank-one flow matrix gives
\[
      x(\zeta)
        =X.\mathrm{flow}\bigl(u,(\zeta/n_0^2)_\tau,x_0\bigr).
 \tag{HM2}\label{eq:sr-source-flow-rescale}
\]
The source expansion is the convergent sum
\[
 P(x(\zeta))=\sum_\nu a_\nu
       \exp\!\left({4\pi i\zeta\over n_0^2}
                              \sum_\tau\nu_\tau\right)
                  \qquad(\operatorname{Im}\zeta\ge0),
 \tag{HM3}\label{eq:sr-source-isotropic-expansion}
\]
where \(a_\nu=X.\mathrm{fj}(\lambda,u,\nu,P,x_0)\).
Every nonzero \(a_\nu\) has all \(\nu_\tau\ge0\), and
\(a_0=0\) by the defining zero-coefficient condition of \(SS\).
Thus every nonzero term has \(\sum_\tau\nu_\tau>0\).

At \(\zeta=0\), the source's \(\mathrm{HasSum}\) in the finite
dimensional complex target gives \(\sum_\nu|a_\nu|<\infty\).
For \(\operatorname{Im}\zeta\ge0\), each summand in (HM3) has
norm at most \(|a_\nu|\), and its norm tends to zero as
\(\operatorname{Im}\zeta\to+\infty\).
Dominated convergence for this summable family therefore yields
\[
                         P(x(\zeta))\longrightarrow0 .
 \tag{HM4}\label{eq:sr-source-flow-decay}
\]
The argument applies to the source's function equal to
\(P(x(\zeta))\) on the flow domain and zero elsewhere, since its
values agree for all sufficiently large imaginary parts.
This is
\srcid{MaximumBoundField.flowTermNormDecay.tendsto_zero}.
It uses the actual Fourier expansion and positive-frequency
supplier; it does not assume rapid decay for an arbitrary
holomorphic function.

\subsection{The maximum estimate before host rescaling}
Set \(f=[E:\mathbb Q]/2\) and
\[
 \varepsilon_*=\min\!\left\{1,\,
 {\pi\over M(fw+2)k_0\log2}\right\}.
 \tag{HM5}\label{eq:sr-source-epstar}
\]
Assume \(M,w>0\), \(P\in SS_\lambda(X)\),
\(\sum_\tau\operatorname{spread}_\tau(\lambda)\le w\), and that
the source rational transvections are in the arithmetic group
with trivial character:
for every isotropic \(v\in L'\) and every integral
conjugation-fixed \(b\in E\), there is \(\gamma\in\Gamma_X\)
such that \(\chi_\lambda(\gamma)=1\) and
\[
                  \tau(\gamma)=\mathrm{flowMat}_{v,\tau}
                                                  (M\tau(b)).
 \tag{HM6}\label{eq:sr-source-flow-period}
\]
Let \(x_0\) be a positive half maximum:
\[
 0<\|P\|_{\lambda,x_0},\qquad
 \forall x\in\mathcal D_X:\quad
                 \|P\|_{\lambda,x}\le2\|P\|_{\lambda,x_0}.
 \tag{HM7}\label{eq:sr-source-half-max}
\]

\begin{srproposition}\label{prop:sr-source-maxbound}
For every \(0\ne v\in L'\), there is a good embedding \(\tau\)
with
\[
                 \varepsilon_*
                  \le\operatorname{maj}_{x_0,\tau}(\tau(v)).
 \tag{HM8}\label{eq:sr-source-maxbound}
\]
\end{srproposition}
\begin{proof}
Suppose every displayed majorant \(m_\tau\) is
less than \(\varepsilon_*\le1\).
The value \(\langle v,v\rangle\) is integral, and its every
complex conjugate has absolute value at most \(m_\tau<1\);
the other half of the embeddings are complex conjugates.
Its rational norm is therefore an integer of absolute value
less than one, so it vanishes.
The source majorant inequality thus makes \(v\) isotropic.

Each \(m_\tau>0\), since \(v\ne0\).
The explicit source flow-domain criterion is
\[
 x(\zeta)\in\mathcal D_X
      \quad\Longleftrightarrow\quad
                 1+k_0m_\tau\operatorname{Im}\zeta>0
                       \quad\text{for every }\tau .
 \tag{HM9}\label{eq:sr-source-flow-domain}
\]
Put \(T=(k_0\varepsilon_*)^{-1}\) and
\(t=(2k_0\varepsilon_*)^{-1}\), so \(0<t<T\).
The assumed strict inequalities imply that the flow lies
in the domain on \(\operatorname{Im}\zeta>-T\).
The flow matrix is polynomial in \(\zeta\), so
\(g(\zeta)=P(x(\zeta))\) is holomorphic there.
Taking \(b=1\) in (HM6), using the flow composition law and
the section's arithmetic transformation, makes \(g\)
\(M\)-periodic.
Equation (HM4) supplies its limit zero at imaginary infinity.

On the line \(\operatorname{Im}\zeta=-t\), all
\(1+k_0m_\tau\operatorname{Im}\zeta\) lie in \((1/2,1]\).
The source Gram-minor comparison for the isotropic flow gives
\[
 {n_\lambda(x_0)\over n_\lambda(x(\zeta))}
 \le\prod_\tau
   \max\{1+k_0m_\tau\operatorname{Im}\zeta,\,
             (1+k_0m_\tau\operatorname{Im}\zeta)^{-1}\}
                           ^{\operatorname{spread}_\tau(\lambda)/2}
 \le2^{w/2}.
\]
Together with (HM7), this implies
\[
     |g(\zeta)|\le
           {2\,\|P\|_{\lambda,x_0}\,2^{w/2}\over n_\lambda(x_0)}.
\]
Apply (HM1) and multiply by \(n_\lambda(x_0)>0\).
Since \(g(0)=P(x_0)\) and the half-maximum norm is positive,
\[
                  1\le e^{-2\pi t/M}\,2\,2^{w/2}.
 \tag{HM10}\label{eq:sr-source-maxbound-contradiction}
\]
But (HM5) gives
\(2\pi t/M=\pi/(Mk_0\varepsilon_*)\ge(fw+2)\log2\).
Since \(f\ge1\) and \(w>0\), the logarithm of the right side
of (HM10) is at most
\(-(fw+2)\log2+(1+w/2)\log2<0\).
This contradicts (HM10).
The proof is exactly \srcid{MaxBoundPf.maxBound_holds}.
The flow-domain and Gram-minor comparisons invoked above are the
source \srcid{FlowDomainCriterionHolds} and \srcid{M3bNormHolds}.
\end{proof}

\subsection{The host scalar and the root discriminant}
Write \(d_E=|\operatorname{disc}E|\) and
\(\mathrm{rd}_E=d_E^{1/[E:\mathbb Q]}\).
The principal host's two fractional ideals \(A,\bar A\) satisfy
\[
                 (\xi_0)A\bar A\,\mathfrak D_E=O_E,\qquad
                 |\tau(\xi_0)|=x>0\quad\text{for every }\tau .
 \tag{HM11}\label{eq:sr-source-host-ideal}
\]
Here \(\mathfrak D_E\) is the different.

\begin{srlemma}\label{lem:sr-source-small-inverse}
There is \(b\in E\), \(b\ne0\), such that
\(bA\subset O_E\) and
\[
                        |\tau(b)|^2\le x\,\mathrm{rd}_E^2
                            \quad\text{for every }\tau.
 \tag{HM12}\label{eq:sr-source-small-inverse}
\]
\end{srlemma}
\begin{proof}
First suppose \(N(\bar A)\le N(A)\).
Taking norms in (HM11), with \([E:\mathbb Q]=2f\), gives
\[
                  x^{2f}N(A)N(\bar A)d_E=1,
       \qquad N(A)^{-1}\le x^f\sqrt{d_E}.
\]
Set \(R=\sqrt{x\,\mathrm{rd}_E^2}\).
For the fractional ideal \(A^{-1}\), the complex-disc version
of Minkowski's theorem applies with this same radius at every
infinite place, because
\[
 2^fN(A)^{-1}\sqrt{d_E}
        \le2^fx^fd_E
        <\pi^fx^fd_E=\pi^fR^{2f}.
\]
It supplies \(0\ne b\in A^{-1}\) with every \(|\tau(b)|<R\).
This gives (HM12) and \(bA\subset O_E\).
If the norms have the opposite order, apply this argument with
\(A,\bar A\) reversed and conjugate \(b\).
The product identity is unchanged, conjugation exchanges the
ideals and permutes all embeddings, giving the same conclusion.
These are the source
\srcid{FirstMinAtHalfMax.exists_small_inv} and
\srcid{exists_small_inv_AA}.
\end{proof}

\subsection{The actual first minimum and the orthogonal half maximum}
Now fix the actual host packet \(D\), stage \(j\), supplied
uniformization laws \(\mathrm{Laws}(D,j,X_o,c,U,\mathrm{dict})\),
a good embedding \(\sigma\), and one component \(k\).
The source's principal-level construction gives a datum \(X\)
with this same underlying component, together with
\[
 \Lambda=A L',\quad k_0=2\,\xi_R(D),\quad
 M'=l.M\,N(\mathrm{ccK})\,D.\mathrm{host}.d,\quad
 M'k_0=2\,l.M\sqrt{D.\mathrm{host}.d},
 \tag{HM13}\label{eq:sr-source-transvection-data}
\]
where \(\Lambda=AL'\) means the actual \(O_E\)-span of
\(au\), \(a\in A,u\in L'\).
The transvections in (HM6) hold with \(M'\).
This precise supplier is
\srcid{ExistsSelfDualDataTransvectionPrincipalLevel.exists_dataX_transv}.
It uses the actual perfect trace form and the actual principal
congruence inclusion, not the abstract lattice alone.

Define
\[
 \epsilon_{\rm lev}(w)=
 \min\!\left\{1,\,
 {\pi\over2\,l.M\sqrt{D.\mathrm{host}.d}\,
                    (fw+2)\log2}\right\}.
 \tag{HM14}\label{eq:sr-source-epslev}
\]
Equation (HM13) identifies this with (HM5).
The source's \(\operatorname{firstMinimumSq}\) is the infimum,
over nonzero \(v\in\Lambda\), of
\(\sum_\tau 2\xi_R(D)\operatorname{maj}_{x_0,\tau}(\tau(v))\).

\begin{srproposition}\label{prop:sr-source-first-minimum-halfmax}
At a positive half maximum of a section with weight spread at
most \(w>0\) and trivial arithmetic character,
\[
 \operatorname{firstMinimumSq}(x_0)
             \ge {2\,\epsilon_{\rm lev}(w)\over\mathrm{rd}_E^2}.
 \tag{HM15}\label{eq:sr-source-first-minimum}
\]
\end{srproposition}
\begin{proof}
Choose \(b\) from (HM12) with \(x=\xi_R(D)\).
For \(v\in\Lambda\), the exact span in (HM13) and
\(bA\subset O_E\) give \(bv\in L'\).
For \(v\ne0\), (HM8) therefore supplies \(\tau_0\) with
\[
 \epsilon_{\rm lev}(w)
       \le\operatorname{maj}_{x_0,\tau_0}(\tau_0(bv))
       =|\tau_0(b)|^2\operatorname{maj}_{x_0,\tau_0}(\tau_0(v))
       \le\xi_R(D)\mathrm{rd}_E^2
                       \operatorname{maj}_{x_0,\tau_0}(\tau_0(v)).
\]
Every term in the first-minimum sum is nonnegative.
Keeping its \(\tau_0\)-term proves the claimed lower bound for
every nonzero \(v\).
The lattice spans \(E^r\), and the actual packet has \(r\ge1\),
so the set defining the infimum is nonempty.
Taking the infimum proves (HM15).
This is \srcid{FirstMinAtHalfMax.firstMin_at_halfMax}.
\end{proof}

\begin{srproposition}\label{prop:sr-source-orthogonal-halfmax}
Suppose a family \(P\) belongs to the actual
\(\mathrm{holSub}(U_\sigma,\lambda)\), has positive finite
global supremum \(S\), has trivial arithmetic characters, and
every component has spread at most the same \(w>0\).
There are a component \(k\) and an orthogonal frame \(x_f\) in
its domain such that
\[
              S\le2|P_k(x_f)|,\qquad
 \operatorname{firstMinimumSq}(x_f)
             \ge {2\,\epsilon_{\rm lev}(w)\over\mathrm{rd}_E^2}.
 \tag{HM16}\label{eq:sr-source-orthogonal-halfmax}
\]
\end{srproposition}
\begin{proof}
Since \(0<S<\infty\), the supremum definition gives a frame
\(x\) in one component with \(S/2<\|P_k\|_{\lambda,x}\).
All component norms are at most \(S\), so \(x\) is a positive
half maximum in that component.
The component belongs to \(SS\), because the source
\(\mathrm{holSub}\) is the span of such families and \(SS\)
is a submodule.
Apply (HM15).
Use the actual upper triangular orthogonalizing factor \(b(x)\)
from \srcid{FrameDegreeData.bOf_mem}.
Right multiplication by it preserves the frame domain, the
weighted norm, and \(\operatorname{firstMinimumSq}\).
At the resulting orthogonal frame all successive Gram-minor
norms are one; its weighted norm is therefore \(|P_k(x_f)|\).
This gives both inequalities in (HM16).
It is \srcid{FirstMinAtHalfMax.half_sup_at_ortho}.
No maximum-attainment or compactness assertion is used.
\end{proof}

\paragraph{Frozen-source ledger and scope.}
All sources are read at \srcid{alpoge-source-7cbf7a0c}.
The periodic Schwarz proof is
\srcid{src/ReciprocityTheory/ReductiveGroups/ArithmeticGeometry/FundamentalDomains/PeriodicSchwarzLemma.lean},
blob \srcid{0f49e1e5cf5bd6b20b73373433d9600b5f02d11f}.
In that same directory,
\srcid{MaximumBoundField.lean} is
\srcid{c7b65dca30b0f3b9a858814b7e5bf9e5f42bb774},
\srcid{MaxBoundPf.lean} is
\srcid{dcd344008b85a949193a31cd613956653dddeb08}, and
\srcid{CuspidalExponentialDecay.lean} is
\srcid{a05dc166cf19232be5e2c25ba52b9a7713d53d3e}.
The host-scalar, first-minimum and half-sup proof is
\srcid{src/Development/AbelianVarietyTheory/Proofs/DivisorProofs/FirstMinAtHalfMax.lean},
blob \srcid{63e29ac9732ebfa986c8251e1b2f6917e907e707}.
The exact principal-level supplier (HM13) is
\srcid{ExistsSelfDualDataTransvectionPrincipalLevel.lean} in
that directory, blob \srcid{8a14c6e08bbc4bcc1218f16298c4ed44e1b165d4}.

The source Fourier expansion \srcid{M3bExpHolds}, coefficient
positivity, flow-domain formula \srcid{FlowDomainCriterionHolds},
and Gram-minor comparison \srcid{M3bNormHolds} are retained as
their exact analytic suppliers in the foregoing proof.
The actual perfect trace-form and principal-congruence
construction in (HM13), and the same uniformization's Borel
orthogonalization identities in (HM16), are likewise the
specified source inputs. The pointwise theorem does not assert
a bound uniform over independently chosen metrics or frames.
